\chapter{Fase 3: DOSBox-Staging y Hugo}

\section{Objetivo}

Tener Hugo ejecutándose en la Pi, mostrado en pantalla por HDMI y jugable con teclado USB. Esta es la última pieza individual antes de unirlo todo en la Fase 4.

\section{Estado legal del juego}

Hugo es \emph{abandonware}: software comercial cuyo soporte por parte del editor original cesó hace décadas y cuya distribución pasiva es ampliamente tolerada aunque no formalmente autorizada. Sitios como \texttt{myabandonware.com} y \texttt{abandonwaregames.net} alojan copias descargables. La situación legal varía por jurisdicción, pero el uso doméstico y educativo no ha sido objeto de litigios.

\section{Paso 1: Descargar Hugo}

Desde la Mac, abrir en un navegador:

\begin{plaincode}
https://www.myabandonware.com/game/hugo-eya
\end{plaincode}

Descargar el archivo ZIP que ofrece el sitio. Una vez en \texttt{\~{}/Downloads/}, transferirlo a la Pi:

\begin{shellcode}
scp ~/Downloads/Hugo*.zip beto@hugo-pbx.local:/tmp/
\end{shellcode}

En la Pi, descomprimir al directorio del juego:

\begin{shellcode}
sudo mkdir -p /opt/hugo-phone/games/hugo
sudo unzip /tmp/Hugo*.zip -d /opt/hugo-phone/games/hugo
sudo chown -R beto:beto /opt/hugo-phone/games/hugo
ls /opt/hugo-phone/games/hugo/
\end{shellcode}

Debería listarse \texttt{hugo.exe} (o un archivo similar). Algunas distribuciones traen un instalador \texttt{install.exe} que se debe ejecutar dentro de DOSBox la primera vez.

\section{Paso 2: Verificar configuración de DOSBox-Staging}

El script de instalación copió un \texttt{dosbox-staging.conf} optimizado para Hugo en la Pi. Verificar:

\begin{shellcode}
cat ~/.config/dosbox/dosbox-staging.conf | head -30
\end{shellcode}

Si el ejecutable del juego no se llama exactamente \texttt{hugo.exe}, editar la sección \texttt{[autoexec]}:

\begin{shellcode}
nano ~/.config/dosbox/dosbox-staging.conf
\end{shellcode}

La sección a ajustar es:

\begin{plaincode}
[autoexec]
mount c /opt/hugo-phone/games
c:
cd hugo
hugo.exe
exit
\end{plaincode}

\section{Paso 3: Modos de salida gráfica}

DOSBox-Staging puede dibujar de varias formas. En orden de simplicidad creciente:

\subsection{Modo desarrollo: con escritorio ligero}

Para pruebas y diagnóstico, instalar Openbox (un gestor de ventanas minimalista):

\begin{shellcode}
sudo apt install -y xinit openbox
echo "exec openbox-session" > ~/.xinitrc
startx /usr/games/dosbox-staging
\end{shellcode}

Esto inicia un servidor X11 con Openbox y lanza DOSBox-Staging dentro. Útil para experimentar.

\subsection{Modo appliance: KMSDRM directo}

Para la operación final, DOSBox-Staging puede dibujar directamente al framebuffer del kernel sin servidor X, lo cual reduce uso de memoria y elimina latencia de composición:

\begin{shellcode}
SDL_VIDEODRIVER=kmsdrm dosbox-staging
\end{shellcode}

\begin{nota}
KMSDRM solo funciona desde una TTY local (no por SSH X-forwarding). Es el modo recomendado para la operación normal del appliance.
\end{nota}

\section{Paso 4: Arranque automático del juego}

Configurar la Pi para que al encender entre directamente a Hugo, sin intervención del usuario.

\subsection{Habilitar autologin en TTY1}

\begin{shellcode}
sudo raspi-config
\end{shellcode}

Navegar a: \textsf{System Options} $\to$ \textsf{Boot / Auto Login} $\to$ \textsf{Console Autologin}. Guardar y salir.

\subsection{Ejecutar DOSBox al iniciar sesión}

Editar el \texttt{.bashrc} del usuario \texttt{beto}:

\begin{shellcode}
nano ~/.bashrc
\end{shellcode}

Agregar al final:

\begin{plaincode}
if [ -z "$DISPLAY" ] && [ "$(tty)" = "/dev/tty1" ]; then
  SDL_VIDEODRIVER=kmsdrm dosbox-staging
fi
\end{plaincode}

La condición evita que DOSBox se ejecute si se ingresa al sistema por SSH o desde una TTY distinta de la 1.

Al reiniciar, la secuencia debe ser:

\begin{enumerate}
  \item Boot del kernel
  \item Autologin como \texttt{beto} en \texttt{tty1}
  \item Ejecución de DOSBox-Staging en modo KMSDRM
  \item Autoejecución de Hugo según \texttt{[autoexec]}
\end{enumerate}

\section{Paso 5: Prueba de jugabilidad}

Con un teclado USB conectado a la Pi y Hugo cargado, probar:

\begin{itemize}
  \item Flechas direccionales $\to$ movimiento del personaje
  \item Barra espaciadora $\to$ acción
  \item \texttt{Esc} $\to$ menú/salir
\end{itemize}

Si Hugo se siente excesivamente rápido o lento, ajustar la línea \texttt{cycles} en \texttt{dosbox-staging.conf}:

\begin{tomlcode}
[cpu]
cycles = fixed 8000   # Aumentar si va lento, disminuir si va rapido
\end{tomlcode}

Hugo original requería aproximadamente un 386SX a 25\,MHz. Empezar con \texttt{cycles = fixed 8000} y calibrar.

\section{Paso 6: Salir limpiamente}

Dos pulsaciones de \texttt{Esc} salen al menú de Hugo. La opción \emph{Quit} regresa al prompt de DOSBox; el comando \texttt{exit} del bloque \texttt{[autoexec]} cierra DOSBox automáticamente.

Para salir de emergencia (por ejemplo si el juego se cuelga): \texttt{Ctrl}+\texttt{F9} es el \emph{kill switch} de DOSBox-Staging.

\section{Resolución de problemas}

\subsection{Hugo se cuelga en la intro o las cinemáticas}

DOSBox-Staging en arquitectura ARM ocasionalmente tiene problemas con la decodificación de animaciones VGA. Pulsar cualquier tecla para saltar la intro suele resolverlo. Si persiste, bajar la resolución de salida:

\begin{tomlcode}
[sdl]
windowresolution = 640x480
\end{tomlcode}

\subsection{El audio se entrecorta}

Aumentar el buffer del mezclador:

\begin{tomlcode}
[mixer]
prebuffer = 50
blocksize = 2048
\end{tomlcode}

Si aún hay cortes, bajar la calidad: \texttt{rate = 22050}.

\subsection{Pantalla negra al iniciar}

Probablemente el driver KMSDRM no está disponible en la versión instalada de SDL2. Alternativas:

\begin{shellcode}
SDL_VIDEODRIVER=fbcon dosbox-staging
\end{shellcode}

O instalar el escritorio mínimo:

\begin{shellcode}
sudo apt install -y xinit openbox
\end{shellcode}

Y lanzar DOSBox dentro de \texttt{startx}.

\subsection{Hugo arranca pero no responde al teclado}

\begin{itemize}
  \item Verificar que \texttt{core = dynamic} esté presente en la sección \texttt{[cpu]}.
  \item Algunos títulos requieren \texttt{cputype = 386} en lugar de \texttt{486\_slow}.
\end{itemize}

\section{Cierre de fase}

En este punto el juego corre con teclado USB. La siguiente fase reemplaza ese teclado USB por uno virtual controlado desde el teléfono.
